% =============================================================================
%  扉页、摘要、译者说明
% =============================================================================
\begin{titlepage}
\centering
\vspace*{3.2cm}
{\Huge\bfseries Arkhipov–Bezrukavnikov 等价\par}
\vspace{0.9em}
{\large\itshape The Arkhipov–Bezrukavnikov Equivalence\par}
\vspace{2.4em}
{\large Colton Sandvik\par}
\vspace{0.4em}
{\normalsize BunG 讨论班报告（2025 年 1 月 29 日）\par}
\vspace{3.0em}
{\large 中译注本\par}
\vspace{0.6em}
{\normalsize 数译 · 几何 Langlands 系列 · SY02-40\par}
\vfill
{\small
\begin{tabular}{@{}rl@{}}
底本 & Colton Sandvik, \emph{The Arkhipov–Bezrukavnikov Equivalence},\\
     & BunG Seminar, 2025-01-29，13 页（pdfTeX 排版）\\
译者 & 数译（DSH agent：lead）\\
版本 & v1.0，2026 年 11 月 26 日\\
体例 & 公式编号 (1)–(8)、定理编号「节.小节.序号」均照原文；\\
     & 译注（橙色框）与「本节重点」（蓝色框）为译者所加\\
\end{tabular}}
\end{titlepage}

\section*{摘要}
\addcontentsline{toc}{section}{摘要}

本讲义是作者在 BunG 讨论班上关于 Arkhipov–Bezrukavnikov 等价所作报告的配套材料。
我们首先搭建中心层（central sheaf）与 Wakimoto 层的机器。
然后，我们解释 Springer 解消上凝聚层（coherent sheaf）的几何。
我们以此构造一个从 Springer 解消上的凝聚层到反球 Hecke 范畴（antispherical Hecke category）的函子。
最后，我们把反球 Hecke 范畴与 Iwahori–Whittaker 范畴联系起来。

\section*{译者说明}
\addcontentsline{toc}{section}{译者说明}

\subsection*{一、这份译本的定位}

底本是 Colton Sandvik 于 2025 年 1 月 29 日在 BunG 讨论班上的报告讲义，13 页，未正式发表。
它用大约五节的篇幅，把 Arkhipov–Bezrukavnikov 等价
\[
D^b\Coh^{G^\vee}(\tN)\;\simeq\;D^b_{\IW}(\Fl_G)
\]
的证明路线完整地走了一遍：从 Gaitsgory 中心函子与 Wakimoto 层出发，经 Springer 解消上的凝聚层，
落到反球 Hecke 范畴与 Iwahori–Whittaker 层，最后用「本质满 + 全忠实」两半合拢。
全文没有一句废话，每一节的每个构造都在后面被用到，适合作为进入 \mbox{[ArBe]} 原文的导引。

译文以「读懂原文」为第一目标，因此做了三件事：
（1）公式、编号、定理编号与原文逐一对齐，便于左右对读；
（2）为原文中因篇幅而省略、或读者容易误解的地方补写译注（橙色框），共 26 条；
（3）每节末给出「本节重点」（蓝色框），点明该节在整个证明链条中的位置。

\subsection*{二、体例}

\begin{enumerate}[leftmargin=2.2em,itemsep=2pt]
\item \textbf{编号}：全文只有 8 个编号公式，即 (1)–(8)，与原文一致；
定理、命题、引理、推论、定义、注、旁注共用一个计数器，编号形如 1.0.1、2.2.1、5.0.1
（§1、§5 之下没有小节，故中段为 0）。引用时一律沿用原编号。
\item \textbf{术语}：首次出现写作「中文（English）」；此后只用中文。
术语表见 \texttt{工程文件/GLOSSARY.tsv}。
\item \textbf{译注}：只写三类内容——术语背景、原文含糊或省略、原文笔误；每条 1–4 句，不代替读者作证明。
原文笔误一律\textbf{照录}，并在译注中以「[原文疑误]」标出正确写法，不改动原文的数学内容。
\item \textbf{公式}：全部依据底本页面图像录入（150 dpi 整页图与局部 450–600 dpi 放大图），
不依赖 PDF 文本层抽取，因为后者会丢失上下标与正负号。
\end{enumerate}

\subsection*{三、原文疑误一览}

下列各处按原文照录，并在正文相应位置以译注标出。列在这里，是为了让读者一眼看到「哪些地方不能照抄」。

\begin{itemize}[leftmargin=2.2em,itemsep=3pt]
\item \textbf{p.3}：$\Gr_G^{\Cen,\times}:=f^{-1}(\bbC^\times)=\Fl_G$（第二次出现）——
      与上一式逐字相同，显然是笔误；应为 $\Gr_G^{\Cen,0}:=f^{-1}(0)=\Fl_G$。
\item \textbf{p.10}：$\bigl(Z^{\IW}(W):\Delta^{\IW}_\lambda\bigr)=(\cdots)$——
      同一式中另一侧写作 $Z^{\IW}(V)$，此处的 $W$ 应为 $V$。
\item \textbf{p.11}：图 (7) 中列 $Z(\Sat(L_\lambda))$——
      该行左列为 $V\otimes\cO_{\tN}$、右列为 $\mathcal{T}^{\IW}(V)$，中列应为 $Z(\Sat(V))$。
\item \textbf{p.4}：there are canonical \emph{isomorphsisms}——应为 isomorphisms。
\item \textbf{p.7}：we replace $\PWak$ \emph{be} a (non-full) subcategory——应为 by。
\item \textbf{p.9}：the pro-unipotent radical \emph{if} $I^+$——应为 of。
\item \textbf{p.10}：\emph{There} restriction of $\AvIW$——应为 The。
\item \textbf{p.12}：\emph{There} functor $F_{\IW}$——应为 The。
\item \textbf{p.12 脚注 11}：the Krillov model——应为 Kirillov。
\end{itemize}

\subsection*{四、配套材料}

\begin{itemize}[leftmargin=2.2em,itemsep=2pt]
\item \textbf{资料卡}：\texttt{datasets/ab\_equiv/}（24 张，JSON；
     索引 \texttt{index.jsonl}），
     另有排印本 \texttt{成品/AB\_资料卡\_v1.0.pdf}，可单独打印。
\item \textbf{核对表}：\texttt{译对/} 下有结构核查表与 8 个编号公式、14 条脚注的逐条对照表。
\item \textbf{术语表}：\texttt{工程文件/GLOSSARY.tsv}（约 90 条，含「易混淆点」列）。
\end{itemize}

\subsection*{五、读法建议}

原文只有 13 页，却串起了五节各自独立的技术。第一次读，建议按下面的次序，而不是从头到尾：

\begin{enumerate}[leftmargin=2.2em,itemsep=2pt]
\item 先读 \S\ref{sec:padic}（$p$ 进动机）与定理 \ref{thm:Z}（中心函子的五条性质），
      抓住「中心层 = 把球侧的层搬到旗簇侧、并且落在中心里」这唯一的直觉；
      \S\ref{sec:gr} 与 \S\ref{sec:central} 是为了说清这句话的技术准备。
\item 再读 \S\ref{sec:springer} 的图 (4) 与命题 \ref{prop:Kquot}。
      这张图是全篇几何上唯一「不好画」的地方：$\tN$、$\hN$、$\hNX$ 三个空间与两条竖直的闭浸入，
      记住「$\hN$ 是 $\tN$ 的 $T^\vee$-主丛、$\hNX$ 是它的完备化」就够了。
\item 然后读定理 \ref{thm:av} 与定理 \ref{thm:tilt}，即第 4 节的两条主定理。
      它们的证明都是「梗概」，原文只给思路：一条靠轨道考察与 $\mathrm{IC}_w$ 的消没，
      一条靠对 $G$ 与 $V$ 的双重归约。
\item 最后读第 5 节。此时图 (7) 应当已经不需要解释，而定理 \ref{thm:main} 的证明只剩「五引理」这一步。
\end{enumerate}

标题为「术语背景」的译注可以第一遍跳过；标着「[原文疑误]」的四条译注
（分别在 §\ref{sec:central} 的第二个等式、注 4.0.1、定理 \ref{thm:tilt}、图 (7)）
建议对照原文页面看一遍，因为它们是原文里真正会让人卡住的地方。
